\chapter{Section-Space Histories and Base Averaging}
\label{ch:section-space-histories}

This chapter asks what becomes of the inference histories and Fisher clocks of
\cref{ch:relational-inference} when the state of an agent is a whole section over
a contextual domain of positive dimension, rather than a point of its fibers
over \(c_*\). Three results carry the answer. A free, proper, isometric gauge
action on a strong Riemannian Hilbert manifold of configurations has a smooth
principal quotient on which the gauge-quotient speed is attained by a unique
orthogonal lift, so quotient durations are well defined
(\cref{thm:hist-strong-hilbert-quotient}). The same conclusion fails for a
weak Fisher metric even when the smooth quotient and a topological splitting
exist, and the trace quotient of that example is not a fixed observation
channel (\cref{prop:hist-proper-weak-fisher-quotient,prop:hist-weak-trace-not-markov}).
Averaging sections over the fibers of a base map loses integrated Fisher
information by an exact channel, weight, and context term under a joint-convexity
condition on the fiber chart, and without that condition the loss can be
negative (\cref{thm:hist-averaging-defect,prop:hist-averaging-convexity}). The
Hilbert theorem is the section-space counterpart of
\cref{prop:hist-invariant-history-horizontal}, and the averaging defect is the
counterpart of \cref{thm:hist-pointwise-contraction-lift} for a base map that
collapses contexts; the covariant first jets and base-direction pullbacks used
here are those of \cref{ch:pullback-geometry}.

Over the single point \(c_*\) every construction of the chapter reduces to a
fiber statement of the main manuscript. A section is then a point of the
fiber, every base curve is constant, so every nonstationary curve is vertical
and horizontality is vacuous; the design and continuum clocks reduce to the
pointwise clock \eqref{eq:hist-pointwise-clock-speed}; the Hilbert quotient
becomes one orthogonal projection; and averaging over the fibers of the identity
of \(\{c_*\}\) is the pointwise map of \cref{thm:hist-pointwise-contraction-lift}.

The chapter declines to claim that the set of fine sections on which a
pointwise bundle morphism induces a configuration map (\cref{thm:pb-section-descent})
is a smooth submanifold carrying a smooth induced arrow, so the coarse
configuration map of \cref{sec:hist-base-averaging} is declared rather than
constructed by section descent; that the base, the principal bundle, or the
free energy selects a base measure, channel weights, or a strong manifold
topology for a continuum clock; that the Hilbert theorem supplies a singular
quotient when the action is not free or not proper; and that the tiers
exhibited here classify section-space metrics. \status{NOT-CLAIMED} Whether a
declared recognition family can be brought to an admitted configuration tier is
\cref{open:hist-section-tier-admission}; the smooth structure of the projectable
set is \cref{open:hist-projectable-set}; a probabilistic theory on continuum
section spaces is \cref{open:prob-continuum-theory}.

\section{Fixed-base kinematics}
\label{sec:hist-fixed-base-kinematics}

\definitionheading{Vertical, horizontal, and mixed curves}{def:hist-curve-types}
For agent \(i\) with contextual domain \(\mathcal C_i\), restrict the belief and
model bundles to \(\mathcal C_i\) and form
\begin{equation}
\mathcal E_i
=\mathcal E_b|_{\mathcal C_i}
\times_{\mathcal C_i}
\mathcal E_m|_{\mathcal C_i},
\qquad
\varpi_i:\mathcal E_i\longrightarrow\mathcal C_i.
\label{eq:hist-agent-product-bundle}
\end{equation}
The pair \(\boldsymbol\omega_i=(\omega_b,\omega_m)\) of principal connections
supplies a product horizontal distribution with horizontal lift
\(H^{\boldsymbol\omega_i}\) and a vertical projection
\(\operatorname{ver}^{\boldsymbol\omega_i}\). The vertical Fisher form of
\(\mathcal E_i\) is fixed only after positive channel weights or a joint
statistical metric are declared. Let \(\mathscr I\) be a connected interval and
\(\Gamma:\mathscr I\to\mathcal E_i\) a \(C^1\) curve with base curve
\(\gamma=\varpi_i\circ\Gamma\). Relative to \(\boldsymbol\omega_i\) its velocity
splits uniquely as
\begin{equation}
\dot\Gamma
=H^{\boldsymbol\omega_i}_{\Gamma}\dot\gamma
+v_{\boldsymbol\omega_i}(\Gamma),
\qquad
v_{\boldsymbol\omega_i}(\Gamma)
=\operatorname{ver}^{\boldsymbol\omega_i}(\dot\Gamma),
\label{eq:hist-total-curve-splitting}
\end{equation}
and applying \(T\varpi_i\), which is the identity on horizontal lifts and zero
on vertical vectors, recovers \(\dot\gamma\); this is the connection splitting
of \citet[Ch.~II]{Kobayashi1963}. At a parameter value \(r\) the velocity is
\emph{stationary} when \(\dot\Gamma(r)=0\); \emph{vertical} when
\(\dot\gamma(r)=0\) and \(\dot\Gamma(r)\ne0\);
\emph{\(\boldsymbol\omega_i\)-horizontal} when
\(v_{\boldsymbol\omega_i}(\Gamma)(r)=0\) and \(\dot\gamma(r)\ne0\); and
\emph{mixed} when \(\dot\gamma(r)\) and \(v_{\boldsymbol\omega_i}(\Gamma)(r)\)
are both nonzero. These four cases partition the parameter values. A curve is
vertical, horizontal, or mixed on \(\mathscr I\) when it has that type at every
parameter value, so the interval-level labels are not exhaustive. Verticality
depends only on the bundle projection; horizontality and the mixed
decomposition depend on the connections.

A base curve carries no verticality predicate, since the predicate concerns a
total-space velocity, and a curve of section configurations
(\cref{sec:hist-section-curves}) carries one only through its evaluation at a
fixed context, \eqref{eq:hist-pointwise-history-verticality}. A vertical curve
stays in one fiber, whereas a horizontal lift compares points of different
fibers by the chosen connection; neither moves the base.

\propositionheading{Horizontality is connection relative}{prop:hist-horizontal-connection-dependence}
The same nonvertical total-space curve can be mixed for one connection and
horizontal for another. \status{ESTABLISHED}

\noindent\emph{Proof.}
Take the trivial line bundle \(\R_{c}\times\R_{y}\to\R_{c}\) and the curve
\(\Gamma(r)=(r,r)\). For the horizontal distribution spanned by
\(\partial_c\), its vertical velocity is \(\partial_y\), so the curve is mixed.
For the distribution spanned by \(\partial_c+\partial_y\), the same velocity is
horizontal. The base trace \(c=r\) is unchanged. \(\square\)

\propositionheading{Lengths of vertical, horizontal, and section curves}{prop:pb-curve-taxonomy}
Fix one channel \(x\in\{b,m\}\) and write \(\omega=\omega_x\),
\(g^F=g^F_x\), and \(E=\mathcal E_x\). For a \(C^1\) curve
\(\Gamma:\mathscr I\to E\) over \(\gamma\), let
\(v^\omega(\Gamma)=\operatorname{ver}^\omega(\dot\Gamma)\) and let
\(L^\omega(\Gamma)=\int_{\mathscr I}g^F(v^\omega(\Gamma),v^\omega(\Gamma))^{1/2}dr\)
be its vertical Fisher length; on a curve confined to one fiber this is the
Fisher length of \eqref{eq:pb-vertical-fisher-length}. Then \(L^\omega\) is
invariant under orientation-preserving \(C^1\) reparameterization; if
\(\Gamma\) is vertical, \(L^\omega(\Gamma)\) does not depend on the connection;
if \(\Gamma\) is \(\omega\)-horizontal, \(L^\omega(\Gamma)=0\); and for a section
curve \(\Gamma=s\circ\gamma\), with the covariant first jet
\eqref{eq:pb-covariant-first-jet},
\begin{equation}
L^\omega(s\circ\gamma)
=\int_{\mathscr I}
g^F\bigl(D^\omega s(\dot\gamma),D^\omega s(\dot\gamma)\bigr)^{1/2}dr,
\label{eq:pb-section-curve-length}
\end{equation}
which is the length of \(\gamma\) in the Fisher pullback
\eqref{eq:pb-fisher-pullback}. If \(\omega'=\omega+\Delta\omega\), with
\(\Delta\omega\) the descended difference form of
\eqref{eq:pb-connection-difference} and \(\vartheta\) the fundamental vertical
map of \eqref{eq:pb-connection-difference-vertical}, then along every curve
\begin{equation}
v^{\omega'}(\Gamma)
=v^\omega(\Gamma)
+\vartheta_{\Gamma}\bigl(\Delta\omega_{\gamma}(\dot\gamma)\bigr),
\label{eq:pb-curve-connection-change}
\end{equation}
so the length of a mixed curve generally depends on the connection.
\status{ESTABLISHED}

\noindent\emph{Proof.}
Under \(r=\phi(\tilde r)\) the vertical velocity is multiplied by
\(\phi'(\tilde r)>0\), and positive homogeneity of the square root turns that
factor into the Jacobian of the change of variables. For a vertical curve
\(\dot\gamma=0\), so the horizontal part of \(\dot\Gamma\) vanishes for every
connection. Horizontality is the equation \(v^\omega(\Gamma)=0\). For
\(\Gamma=s\circ\gamma\) the chain rule gives
\(\dot\Gamma=Ts(\dot\gamma)\), hence
\(v^\omega(\Gamma)=D^\omega s(\dot\gamma)\), which is
\eqref{eq:pb-section-curve-length}. Finally, replacing \(\omega\) by
\(\omega'\) changes the horizontal lift of \(\dot\gamma\) at \(\Gamma\) by
\(-\vartheta_\Gamma(\Delta\omega_\gamma(\dot\gamma))\) in the sign convention
of \eqref{eq:pb-connection-difference-vertical}, and the vertical part of the
fixed vector \(\dot\Gamma\) changes by the opposite amount; for a section curve
this is \eqref{eq:pb-jet-connection-change}. \(\square\)

\section{Curves of sections and pointwise vertical change}
\label{sec:hist-section-curves}

\hypothesisheading{Regular section configuration}{hyp:hist-regular-section-space}
In this chapter \(\mathfrak S_i\) is a declared space of pairs
\(Q=(q^b,q^m)\) of smooth sections of \(\mathcal E_i\to\mathcal C_i\). Wherever a
tangent-space calculation is used, \(\mathfrak S_i\) is either a
finite-dimensional submanifold of configurations or carries a specified smooth
locally convex manifold structure for which the evaluation
\((Q,c)\mapsto Q(c)\) is differentiable. This excludes the space of all smooth
sections taken without a declared manifold structure, since none follows from
writing the sections down.

Oriented histories in \(\mathfrak S_i\) are the equivalence classes of regular
representatives under orientation-preserving \(C^1\) reparameterization,
\eqref{eq:hist-oriented-reparameterization-equivalence}, exactly as in
\cref{def:hist-oriented-history} with \(\mathfrak S_i\) in place of the
configuration manifold at \(c_*\).

\propositionheading{Section histories change pointwise vertically}{prop:hist-section-history-verticality}
Let \(Q_i=(q^b_i,q^m_i):\mathscr I\to\mathfrak S_i\) be a \(C^1\)
representative and define its adjoint family
\begin{equation}
\widehat{\mathrm{ev}}_i:\mathscr I\times\mathcal C_i\longrightarrow\mathcal E_i,
\qquad
\widehat{\mathrm{ev}}_i(r,c)=Q_i(r)(c),
\qquad
\varpi_i\circ\widehat{\mathrm{ev}}_i(r,c)=c.
\label{eq:hist-adjoint-section-family}
\end{equation}
Then at every fixed context
\begin{equation}
T\varpi_i\bigl(\partial_r\widehat{\mathrm{ev}}_i(r,c)\bigr)=0,
\qquad
\partial_r\widehat{\mathrm{ev}}_i(r,c)
\in V_{Q_i(r)(c)}\mathcal E_i.
\label{eq:hist-pointwise-history-verticality}
\end{equation}
With the connections held fixed, the family of Fisher pullbacks
\eqref{eq:pb-fisher-pullback} of the instantaneous sections,
\begin{equation}
\mathbf P_i(r)
=\Bigl(
\bigl(D^{\omega_b}q^b_i(r)\bigr)^*g^F_b,
\bigl(D^{\omega_m}q^m_i(r)\bigr)^*g^F_m
\Bigr),
\label{eq:hist-perceived-geometry-family}
\end{equation}
is an ordered family of base semimetrics that transforms under
\(\widetilde Q_i=Q_i\circ\phi\) by \(\widetilde{\mathbf P}_i=\mathbf P_i\circ\phi\).
\status{ESTABLISHED}

\noindent\emph{Proof.}
Differentiating \(\varpi_i(\widehat{\mathrm{ev}}_i(r,c))=c\) in \(r\) at fixed
\(c\) gives \eqref{eq:hist-pointwise-history-verticality}. Each entry of
\(\mathbf P_i(r)\) depends on \(r\) only through \(Q_i(r)\), because the
connections are fixed, so reparameterization composes with \(\phi\). \(\square\)

A curve of full sections therefore produces a vertical statistical update at
every fixed context although its state is an entire section, and it orders a
family of perceived base geometries without making the base evolve. A
separately declared connection history would add the correction terms of
\cref{prop:pb-pullback-connection-change} to \eqref{eq:hist-perceived-geometry-family}.

\propositionheading{Diagonal evaluation separates history change from base probing}{prop:hist-diagonal-evaluation}
Let \(Q_i:\mathscr I\to\mathfrak S_i\) be a \(C^1\) representative, let
\(\gamma:\mathscr I\to\mathcal C_i\) be a \(C^1\) base probe, set
\(\Gamma(r)=Q_i(r)(\gamma(r))\), and write
\begin{equation}
D^{\boldsymbol\omega_i}Q_i(r)
=\bigl(D^{\omega_b}q^b_i(r),D^{\omega_m}q^m_i(r)\bigr)
\label{eq:hist-product-covariant-section-jet}
\end{equation}
for the pair of covariant first jets \eqref{eq:pb-covariant-first-jet}. Then
\begin{equation}
v_{\boldsymbol\omega_i}(\Gamma)
=\partial_r\widehat{\mathrm{ev}}_i(r,\gamma(r))
+D^{\boldsymbol\omega_i}Q_i(r)\dot\gamma(r).
\label{eq:hist-diagonal-evaluation-velocity}
\end{equation}
At a parameter value with \(\dot\gamma=0\) the diagonal curve is vertical when
the pointwise history velocity is nonzero and stationary when it vanishes; on
an open subinterval where \(\dot\gamma\ne0\) it is mixed exactly when the sum in
\eqref{eq:hist-diagonal-evaluation-velocity} is nonzero and horizontal when the
sum vanishes. Moreover every smooth section \(Q\) of
\(\mathcal E_i\to\mathcal C_i\) is an embedding; a total-space curve has the
form \(Q\circ\gamma\) exactly when it lies in the image of \(Q\), and then
\(\gamma=\varpi_i\circ\Gamma\) is unique. In particular a nonconstant curve
confined to one fiber is not the image of a base curve under a section.
\status{ESTABLISHED}

\noindent\emph{Proof.}
The chain rule for the jointly \(C^1\) map \((r,c)\mapsto Q_i(r)(c)\), which
\cref{hyp:hist-regular-section-space} supplies, gives
\(\dot\Gamma=\partial_r\widehat{\mathrm{ev}}_i(r,\gamma(r))+TQ_i(r)\dot\gamma\).
The first term is vertical by \eqref{eq:hist-pointwise-history-verticality},
and \(\operatorname{ver}^{\boldsymbol\omega_i}\circ TQ_i(r)=D^{\boldsymbol\omega_i}Q_i(r)\)
by \eqref{eq:pb-covariant-first-jet}; applying the vertical projection gives
\eqref{eq:hist-diagonal-evaluation-velocity}. The type statements follow from
\cref{def:hist-curve-types}, since \(\dot\gamma=0\) makes \(\dot\Gamma\) equal to
its vertical part. Since \(\varpi_i\circ Q=\operatorname{id}\), the section is
injective, \(T\varpi_i\circ TQ=\operatorname{id}\) makes it an immersion, and
\(\varpi_i\) restricted to its image is a continuous inverse, so \(Q\) is a
homeomorphism onto its image. If \(\Gamma=Q\circ\gamma\), then
\(\varpi_i\circ\Gamma=\gamma\); conversely a curve in the image satisfies
\(\Gamma=Q\circ\varpi_i\circ\Gamma\). A curve in one fiber has constant
\(\varpi_i\circ\Gamma\), so its only representation is a constant curve.
\(\square\)

\section{Finite-design and continuum clocks}

The pointwise clock of \eqref{eq:hist-pointwise-clock-speed} uses only the
vertical evaluation velocities of \eqref{eq:hist-pointwise-history-verticality}
at one context. Over a base, two further clocks are available once their data
are declared.

\definitionheading{Finite-design and continuum clocks}{def:hist-design-continuum-clocks}
Fix positive channel weights \(w_b,w_m\) on \(\mathcal C_i\) and a \(C^1\)
section history \(Q_i=(q^b_i,q^m_i)\). For a finite design
\(D=\{c_d\}_{d=1}^M\subset\mathcal C_i\) with positive design weights
\(\mathsf p_d\), the finite-design squared speed is
\begin{equation}
\nu_{i,D}^2
=\sum_{d=1}^M\mathsf p_d\left[
w_b(c_d)g_b^F\bigl(\partial_rq^b_i(c_d),\partial_rq^b_i(c_d)\bigr)
+w_m(c_d)g_m^F\bigl(\partial_rq^m_i(c_d),\partial_rq^m_i(c_d)\bigr)
\right].
\label{eq:hist-finite-design-clock-speed}
\end{equation}
For a finite positive base measure \(\varrho_i\), measurable positive channel
weights, and square-integrable evaluation velocities, the continuum squared
speed is
\begin{equation}
\nu_{i,\varrho}^2
=\int_{\mathcal C_i}\left[
w_b(c)g_b^F\bigl(\partial_rq^b_i(c),\partial_rq^b_i(c)\bigr)
+w_m(c)g_m^F\bigl(\partial_rq^m_i(c),\partial_rq^m_i(c)\bigr)
\right]d\varrho_i(c).
\label{eq:hist-continuum-clock-speed}
\end{equation}
Both clocks are declared data, and without their declarations only the
pointwise clock is defined.

\propositionheading{The design clock has a large radical}{prop:hist-design-clock-radical}
The finite-design speed \eqref{eq:hist-finite-design-clock-speed} is the product
Fisher speed of the declared design and weights, and every section variation
that vanishes at all \(c_d\) lies in its radical. In particular every nonzero
vertical variation in \(T_Q\mathfrak S_i\) supported in an open set disjoint from
\(D\), such as a bump variation when \(\mathfrak S_i\) admits one, is a nonzero
element of the radical, and the design speed is then not positive definite on
\(T_Q\mathfrak S_i\). \status{ESTABLISHED}

\noindent\emph{Proof.}
Each summand of \eqref{eq:hist-finite-design-clock-speed} evaluates the
variation only at a design point. \(\square\)

\lemmaheading{Local continuum metrics see no cross-context pairing}{lem:hist-continuum-clock-locality}
Let a configuration metric on \(T_Q\mathfrak S_i\) have the local single-integral
form obtained by polarizing \eqref{eq:hist-continuum-clock-speed}. Then two
variations whose contextual supports are disjoint have zero pairing.
Consequently, if the exact Fisher bilinear form of a declared joint law of
section configurations pairs two such variations nontrivially, that form is not
of the local single-integral form, and replacing it by
\eqref{eq:hist-continuum-clock-speed} defines a local approximation.
\status{ESTABLISHED}

\noindent\emph{Proof.}
For variations \(V,W\) with disjoint supports the integrand
\(\sum_xw_x(c)g^F_x(V_x(c),W_x(c))\) vanishes at every \(c\). \(\square\)

Contextual dependence of a law alone therefore does not settle whether the
continuum clock is its exact Fisher speed; that identity must be derived from
the score covariance of the declared joint law, and any cross-context term of the
exact metric must be retained.

\section{Strong Hilbert quotients and attained duration}
\label{sec:hist-strong-hilbert-quotient}

At \(c_*\) with finitely many agents the configuration space is
finite-dimensional, the gauge group is a finite-dimensional Lie group, and the
gauge-quotient speed is one orthogonal projection
(\cref{prop:hist-invariant-history-horizontal}). For configurations of sections
the same formula needs a Hilbert category. For a smooth proper action of a
Banach--Lie group on a Banach manifold, the infinitesimal orbit maps have closed
range \citep[Theorem~2.1]{JotzNeeb2008}, and the orbit space of a free proper
action on a Hilbert manifold carries a unique Hilbert manifold structure for
which the quotient map is a submersion \citep[Corollary~2.5]{JotzNeeb2008}; both
use properness in the sense of \eqref{eq:hist-proper-action-map}. The theorem
below adds what a quotient duration needs and the citation does not supply: the
metric-orthogonal splitting, its attainment, the principal bundle structure,
and a strong quotient metric making the projection a Riemannian submersion. Its
proof is given in \cref{sec:hist-strong-hilbert-quotient-proof}.

\theoremheading{Strong Hilbert quotient duration}{thm:hist-strong-hilbert-quotient}
Let \(\mathcal M\) be a nonempty Hausdorff \(C^\infty\) real Hilbert manifold
with a \(C^\infty\) strong Riemannian metric \(\mathsf g\): the metric norm gives
the declared tangent topology and, in each Hilbert chart, its positive metric
operator is boundedly invertible and varies smoothly in operator norm. Let a
Hausdorff \(C^\infty\) real Banach--Lie group \(\mathbb G\), with Lie algebra
\(\mathfrak g\), act jointly smoothly, freely, and isometrically on
\(\mathcal M\), and assume that the action map
\begin{equation}
\mathfrak a:\mathbb G\times\mathcal M\longrightarrow
\mathcal M\times\mathcal M,
\qquad
\mathfrak a(g,Q)=(g\cdot Q,Q),
\label{eq:hist-proper-action-map}
\end{equation}
is closed with compact fibers. Then the infinitesimal action
\begin{equation}
A_Q\xi=\left.\frac{d}{dt}\right|_{t=0}\exp(t\xi)\cdot Q
\label{eq:hist-infinitesimal-orbit-map}
\end{equation}
is injective with closed range \(V_Q\); the spaces \(V_Q\) and
\(H_Q=V_Q^{\perp_{\mathsf g}}\) are smooth complementary subbundles; the orbit
space \(\mathcal M/\mathbb G\) is, componentwise, a Hausdorff smooth Hilbert
manifold; and its projection \(\pi\) is a smooth principal bundle and a
Riemannian submersion for a strong quotient metric \(\overline{\mathsf g}\).
For every \(Q\in\mathcal M\) and \(v\in T_Q\mathcal M\),
\begin{equation}
\|d\pi_Qv\|_{\overline{\mathsf g}}^2
=\min_{\xi\in\mathfrak g}\mathsf g_Q(v-A_Q\xi,v-A_Q\xi),
\label{eq:hist-quotient-gauge-speed}
\end{equation}
the minimizing \(\xi\) is unique and smooth in \((Q,v)\), and the quotient speed
and its integrated duration do not depend on a \(C^1\) history-dependent choice
of representative. Neither metric completeness nor finite dimensionality of
\(\mathbb G\) is assumed. \status{ESTABLISHED}

Applied to a section configuration \((\mathcal M,\mathsf g)=(\mathfrak S_i,\mathsf G_i^F)\),
the theorem requires that the Hilbert topology, the strong metric, the
Banach--Lie gauge group, and the smooth proper isometric action be established
first. Passive re-trivialization changes coordinates and adds no segment to a
history.

\section{A proper action with a weak Fisher quotient of zero speed}
\label{sec:hist-weak-fisher-quotient}

Closed range and a topological complement do not suffice for a weak metric,
even when the smooth quotient and a global topological splitting exist.

\propositionheading{Proper weak-Fisher quotient obstruction}{prop:hist-proper-weak-fisher-quotient}
Let \(\mathcal C=\R/\mathbb Z\) carry normalized arclength \(dc\), and set
\begin{equation}
\mathcal M=H^1(\mathcal C;\R),
\qquad
\|q\|_{H^1}^2=\int_{\mathcal C}(|q|^2+|q'|^2)dc,
\qquad
\mathcal K=\{h\in\mathcal M:h(0)=0\},
\label{eq:hist-weak-sobolev-model}
\end{equation}
using the continuous periodic representative of each Sobolev class. The
additive Hilbert Lie group \(\mathcal K\), with the inherited \(H^1\) topology,
acts by \(h\cdot q=q+h\). Equip \(\mathcal M\) with
\begin{equation}
\mathsf g_q(v,w)=\int_{\mathcal C}v(c)w(c)dc.
\label{eq:hist-weak-integrated-fisher}
\end{equation}
The action is smooth, free, proper, and isometric; its smooth quotient is
\(\R\), via \(q\mapsto q(0)\); and \(\mathcal K\) is closed and complemented in
\(H^1\). Nevertheless
\begin{equation}
\mathcal K^{\perp_{\mathsf g}}=\{0\},
\qquad
\inf_{h\in\mathcal K}\|v-h\|_{\mathsf g}^2=0
\quad\text{for every }v\in H^1,
\label{eq:hist-weak-quotient-collapse}
\end{equation}
and the infimum is attained exactly when \(v(0)=0\). The form \(\mathsf g\) is
the Fisher metric of the identifiable normalized joint observation family
\begin{equation}
P_q(dc,do)=\mathcal N\bigl(o;q(c),1\bigr)(dc\otimes do),
\label{eq:hist-weak-joint-law}
\end{equation}
which is differentiable in quadratic mean in the \(H^1\) parameter topology.
Thus every nonzero quotient tangent has zero induced speed and no minimizing
lift, although the quotient is smooth and proper. \status{ESTABLISHED}

\noindent\emph{Proof.}
Absolute continuity and the Cauchy--Schwarz inequality give
\begin{equation}
\|q\|_\infty\leq\|q\|_2+\|q'\|_2\leq\sqrt2\|q\|_{H^1};
\label{eq:hist-sobolev-trace-bound}
\end{equation}
for the first inequality bound \(|q(c)|\) by \(|q(t)|+\int_{\mathcal C}|q'|\)
and average over \(t\). The trace \(\operatorname{tr}_0(q)=q(0)\) is therefore
bounded, and
\begin{equation}
q\longmapsto\bigl(q-q(0)\mathbf1,q(0)\bigr)
\quad\text{identifies}\quad
H^1\cong\mathcal K\oplus\R
\label{eq:hist-trace-topological-splitting}
\end{equation}
by bounded inverse maps. An \(H^1\)-Cauchy sequence converges uniformly by
\eqref{eq:hist-sobolev-trace-bound} and its derivatives converge in \(L^2\), so
passing to the limit in \(q(c)-q(0)=\int_0^cq'(t)dt\) gives an absolutely
continuous periodic limit with that derivative; this verifies the Hilbert
structure used for the group and the manifold.

The action is translation in the \(\mathcal K\) factor of
\eqref{eq:hist-trace-topological-splitting}; it is smooth and free, and its
quotient is the displayed line. The action map \((h,q)\mapsto(q+h,q)\) is a
homeomorphism onto the closed relation
\(\{(u,q):\operatorname{tr}_0(u-q)=0\}\), with inverse \((u,q)\mapsto(u-q,q)\),
so it is closed with fibers that are singletons or empty, which is properness
in the sense of \eqref{eq:hist-proper-action-map}. Every translation has
identity derivative and preserves \(\mathsf g\).

Let \(\mathsf b(t)=\exp(1-(1-t^2)^{-1})\) for \(|t|<1\) and \(\mathsf b(t)=0\)
otherwise, and for \(n\geq3\) let \(\mathsf b_n(c)=\mathsf b(nc)\) in the signed
circle coordinate \([-1/2,1/2]\), extended periodically. These are smooth with
\(\mathsf b_n(0)=1\), and with \(\mathrm c_1=\int_{-1}^1\mathsf b^2>0\) and
\(\mathrm c_2=\int_{-1}^1(\mathsf b')^2>0\),
\begin{equation}
\|\mathsf b_n\|_2^2=\frac{\mathrm c_1}{n},
\qquad
\|\mathsf b_n'\|_2^2=n\mathrm c_2,
\qquad
v^{(n)}=v-v(0)\mathsf b_n\in\mathcal K.
\label{eq:hist-trace-shrinking-bumps}
\end{equation}
Hence \(\|v-v^{(n)}\|_2=|v(0)|\sqrt{\mathrm c_1/n}\to0\). If \(v\) is
\(\mathsf g\)-orthogonal to \(\mathcal K\), then
\(0=\lim_n\mathsf g(v,v^{(n)})=\|v\|_2^2\), so \(v=0\). The same approximation
gives the zero infimum in \eqref{eq:hist-weak-quotient-collapse}; attainment at
\(h\) would force \(v-h=0\), because the \(L^2\) form is positive definite on
\(H^1\), and this is equivalent to \(v(0)=0\).

The metric is weak in the declared topology: its musical map is injective but
does not represent the continuous trace covector, since
\(\operatorname{tr}_0=\mathsf g(v,\cdot)\) would give
\begin{equation}
1=\operatorname{tr}_0(\mathsf b_n)=\int_{\mathcal C}v\mathsf b_ndc,
\qquad
\left|\int_{\mathcal C}v\mathsf b_ndc\right|
\leq\|v\|_2\sqrt{\mathrm c_1/n}\longrightarrow0.
\label{eq:hist-weak-musical-obstruction}
\end{equation}
The derivative norm in \eqref{eq:hist-trace-shrinking-bumps} diverges, so the
bumps do not converge in \(H^1\), and properness in that topology is not
contradicted.

It remains to show that \eqref{eq:hist-weak-integrated-fisher} is a
statistical tensor. The law \eqref{eq:hist-weak-joint-law} is normalized by
integrating each conditional Gaussian and then the unit circle measure.
Equality of two joint laws gives equality of their conditional means almost
everywhere in \(c\), and continuity gives \(q=\widetilde q\) everywhere, which
is identifiability. Write \(\mathrm r_\mu(o)=\mathcal N(o;\mu,1)^{1/2}\) and
\(\mathrm S(q)(c,o)=\mathrm r_{q(c)}(o)\). The translation-invariant constant
\(\mathrm c_3=\|\partial_\mu^2\mathrm r_\mu\|_{L^2(do)}\) is finite, since the
derivative is a quadratic polynomial times a Gaussian square root. Taylor's
formula in \(L^2(do)\), followed by integration over \(c\), gives
\begin{equation}
\begin{split}
D\mathrm S(q)[v](c,o)&=\tfrac12v(c)\bigl(o-q(c)\bigr)\mathrm S(q)(c,o),\\
\|\mathrm S(q+h)-\mathrm S(q)-D\mathrm S(q)[h]\|_2
&\leq\tfrac12\mathrm c_3\|h\|_4^2
\leq\tfrac12\mathrm c_3\|h\|_\infty\|h\|_2
=O(\|h\|_{H^1}^2).
\end{split}
\label{eq:hist-weak-fisher-dqm}
\end{equation}
This is differentiability in quadratic mean with square-integrable directional
score \(\mathbf s_q[v](c,o)=v(c)(o-q(c))\). The conditional mean and variance
of \(O-q(C)\) are zero and one, so
\begin{equation}
\E_q\bigl[\mathbf s_q[v]\bigr]=0,
\qquad
\E_q\bigl[\mathbf s_q[v]\mathbf s_q[w]\bigr]
=\int_{\mathcal C}v(c)w(c)dc,
\label{eq:hist-weak-fisher-score}
\end{equation}
the Fisher metric of the observed pair \((C,O)\), with \(C\) retained. For each
fixed \(h\in\mathcal K\), the bimeasurable observation change
\((c,o)\mapsto(c,o+h(c))\) pushes \(P_q\) to \(P_{q+h}\), so orbit equivalence
is implemented by a declared coordinate action; it is not equality of these
laws in fixed recorded coordinates. \(\square\)

\propositionheading{The trace quotient is not a fixed observation channel}{prop:hist-weak-trace-not-markov}
No parameter-independent Markov kernel \(K:\mathcal C\times\R\rightsquigarrow\R\)
satisfies
\begin{equation}
P_qK=\mathcal N\bigl(q(0),1\bigr)
\quad\text{for every }q\in H^1(\mathcal C;\R).
\label{eq:hist-trace-channel-impossibility}
\end{equation}
Attaching this normalized output family to the smooth parameter quotient of
\cref{prop:hist-proper-weak-fisher-quotient} therefore gives neither a
statistical coarse channel nor its quotient Fisher metric. \status{ESTABLISHED}

\noindent\emph{Proof.}
Use \(\operatorname{TV}(P,P')=\sup_B|P(B)-P'(B)|\), half the \(L^1\) distance for
densities. Integrating the derivative of the standard normal density gives
\begin{equation}
\int_\R\bigl|\mathcal N(o;t,1)-\mathcal N(o;0,1)\bigr|do
\leq|t|\int_\R\bigl|\partial_o\mathcal N(o;0,1)\bigr|do
=|t|\sqrt{2/\pi},
\qquad
\operatorname{TV}(P_q,P_0)\leq\frac{\|q\|_1}{\sqrt{2\pi}}.
\label{eq:hist-trace-total-variation-bound}
\end{equation}
Fix \(\delta>0\) and take \(q_n=\delta\mathsf b_n\), with the bumps of the
preceding proof. Its trace is \(\delta\), whereas
\(\|q_n\|_1=\delta n^{-1}\int_{-1}^1\mathsf b\to0\). Every Markov kernel
contracts total variation: for an output event \(B\) the function
\(z\mapsto K(z,B)\) takes values in \([0,1]\), and writing
\(K(z,B)=\int_0^1\mathbf1_{\{K(z,B)>t\}}dt\) bounds its expectation difference by
the total variation of the inputs. If \eqref{eq:hist-trace-channel-impossibility}
held, it would give
\(\operatorname{TV}(\mathcal N(\delta,1),\mathcal N(0,1))\to0\). But the fixed
event \((\delta/2,\infty)\) has probability difference
\begin{equation}
\mathcal N(\delta,1)\bigl((\delta/2,\infty)\bigr)
-\mathcal N(0,1)\bigl((\delta/2,\infty)\bigr)
=\int_{-\delta/2}^{\delta/2}\mathcal N(o;0,1)do>0,
\label{eq:hist-trace-output-separation}
\end{equation}
a contradiction. A continuous conditional version at \(C=0\) exists, but
selecting it at this zero-probability context is not the action of a fixed
kernel on one joint observation. \(\square\)

\lemmaheading{Minimizing lifts are orthogonal lifts}{lem:hist-orthogonal-lift}
Let \(V\) be a subspace of an inner-product space \(T\). For \(v\in T\), the
element \(k_*\in V\) minimizes \(\|v-k\|\) over \(k\in V\) exactly when
\(v-k_*\in V^\perp\). Hence every class of \(T/V\) has a minimizing lift exactly
when \(T=V\oplus V^\perp\). \status{ESTABLISHED}

\noindent\emph{Proof.}
Differentiating \(t\mapsto\|v-k_*-tk\|^2\) at \(t=0\) for every \(k\in V\)
proves necessity; if \(v-k_*\in V^\perp\), Pythagoras gives
\(\|v-k\|^2=\|v-k_*\|^2+\|k_*-k\|^2\), which proves sufficiency. A class
\(v+V\) has a minimizing lift exactly when \(v\in V+V^\perp\). \(\square\)

A topological complement is therefore insufficient, as
\cref{prop:hist-proper-weak-fisher-quotient} shows, and a smooth quotient metric
built from minimizing lifts additionally needs a smooth compatible horizontal
structure, which \cref{thm:hist-strong-hilbert-quotient} supplies in the strong
case.

\openproblemheading{Admission of declared recognition families to a configuration tier}{open:hist-section-tier-admission}
The history theorems of \cref{ch:relational-inference} hold on the tiers
admitted by \cref{hyp:hist-standing-configuration-tier}, and
\cref{thm:hist-strong-hilbert-quotient} covers free, proper, isometric actions on
strong Riemannian Hilbert manifolds. A weak metric, a Fr\'echet group or
manifold, a singular action, or a changing support is covered by neither. What
is not settled is whether a recognition family that an application declares, in
particular an infinite-dimensional law fiber or a space of sections, admits at
every scale where a natural-gradient field or a Fisher duration is asserted a
manifold calculus, its actual statistical metric, a smooth quotient, and a
metric-orthogonal splitting. The problem is settled for a declared family by
exhibiting these structures, or by exhibiting a family for which one of them
provably fails while the history theorems are still invoked.

\section{An exhibited finite-dimensional configuration tier}

The finite-dimensional kind of \cref{hyp:hist-standing-configuration-tier}
admits every regular statistical manifold with a declared positive-definite
metric. The following tier is a base-integrated Gaussian instance over a
positive-dimensional base, with a strong constant metric, a rank test for
nondegeneracy, and a closed gauge-orbit tangent, built from bundle data rather
than assumed.

\definitionheading{An exhibited finite-dimensional configuration tier}{def:hist-finite-configuration-tier}
Declare at scale \(\ell\) a compact smooth base \(\mathcal C_\ell\) with a finite
positive Borel measure \(\varrho_\ell\); the trivial principal bundle
\(\mathcal C_\ell\times G\) with \(G=(\R^K,+)\) acting on the fiber by
translation of the mean and \(\omega\) the flat connection; the fiber
\(\mathcal B=\{\mathcal N(\mu,\Sigma_0):\mu\in\R^K\}\) with fixed
\(\Sigma_0\succ0\), whose Fisher form is the constant \(g^F=\Sigma_0^{-1}\); a
measurable weight \(w:\mathcal C_\ell\to(0,\infty)\) bounded above and below; and
smooth fields \(\psi_1,\dots,\psi_n:\mathcal C_\ell\to\R^K\) that are linearly
independent in \(L^2(w\varrho_\ell;\Sigma_0^{-1})\) and whose span contains every
constant field. Set
\begin{equation}
\mathcal Q_\ell
=\left\{s_{\mathbf z}:\mathbf z\in\R^n\right\},
\qquad
s_{\mathbf z}(c)
=\mathcal N\Bigl(\sum_jz_j\psi_j(c),\Sigma_0\Bigr).
\label{eq:hist-finite-configuration-tier}
\end{equation}
Let \(\jmath_\ell:\R^K\to\R^n\) be the unique coefficient map with
\(\sum_j(\jmath_\ell g)_j\psi_j(c)=g\) for every \(c\). Constant mean
translations preserve the span and act on coefficients by
\begin{equation}
\mathscr A^\ell_g(\mathbf z)=\mathbf z+\jmath_\ell g.
\label{eq:hist-finite-tier-coefficient-action}
\end{equation}
Only constant translations are declared: a context-dependent translation
\(c\mapsto g(c)\) maps \(\mathcal Q_\ell\) into itself exactly when it lies in
the span of the \(\psi_j\), so the tier is not closed under general
context-dependent gauge transformations. The configuration metric is the
weighted product of marginal fiber metrics
\begin{equation}
\mathsf G_\ell(V,V)
=\int_{\mathcal C_\ell}w(c)
g^F\bigl(\partial_Vs_{\mathbf z}(c),\partial_Vs_{\mathbf z}(c)\bigr)
d\varrho_\ell(c),
\label{eq:hist-finite-configuration-metric}
\end{equation}
not the pullback of a joint-law Fisher metric.

\theoremheading{The exhibited tier satisfies the regular metric hypothesis}{thm:hist-finite-tier-regularity}
Under \cref{def:hist-finite-configuration-tier}, \(\mathcal Q_\ell\) is a
nonempty smooth manifold diffeomorphic to \(\R^n\), with \(\mathbf z=0\) the
constant configuration \(c\mapsto\mathcal N(0,\Sigma_0)\); the metric
\eqref{eq:hist-finite-configuration-metric} is the constant Gram form
\begin{equation}
\mathsf G_\ell=\mathsf M_\ell,
\qquad
(\mathsf M_\ell)_{jk}
=\int_{\mathcal C_\ell}\psi_j(c)^\top\Sigma_0^{-1}\psi_k(c)w(c)
d\varrho_\ell(c),
\label{eq:hist-configuration-gram}
\end{equation}
independent of \(\mathbf z\), finite, and positive definite, hence strong; every
\(C^1\) functional has the unique gradient \(\mathsf M_\ell^{-1}\nabla\Fenergy\),
and for \(\Fenergy\in C^2\) the natural-gradient field is locally Lipschitz with
locally unique integral curves; and the gauge-orbit tangent of constant mean
translations is the closed complemented subspace \(\jmath_\ell(\R^K)\), so the
minimum in \eqref{eq:hist-quotient-gauge-speed} is attained at the
\(\mathsf M_\ell\)-orthogonal projection. The tier is therefore an instance of
the finite-dimensional kind of \cref{hyp:hist-standing-configuration-tier}.
\status{ESTABLISHED}

\noindent\emph{Proof.}
The bundle is trivial and \(s_{\mathbf z}\) is smooth jointly in
\((c,\mathbf z)\). If \(s_{\mathbf z}=s_{\mathbf y}\), then
\(\sum_j(z_j-y_j)\psi_j=0\) in \(L^2(w\varrho_\ell;\Sigma_0^{-1})\), and
independence gives \(\mathbf z=\mathbf y\), so the family is diffeomorphic to
its coefficient space. Since \(\partial_Vs_{\mathbf z}(c)=\sum_jv_j\psi_j(c)\)
does not depend on \(\mathbf z\), the integrand of
\eqref{eq:hist-finite-configuration-metric} is \(\mathbf v^\top\mathsf M_\ell\mathbf v\);
finiteness follows from compactness of the base and boundedness of \(w\).
If \(\mathbf v^\top\mathsf M_\ell\mathbf v=0\), then \(\sum_jv_j\psi_j=0\) in
\(L^2(w\varrho_\ell;\Sigma_0^{-1})\), so \(\mathbf v=0\). An inner product on a
finite-dimensional space induces its norm topology and has an invertible musical
map, which gives strongness and the gradient. The span condition defines
\(\jmath_\ell\) and gives \(s_{\mathscr A^\ell_g(\mathbf z)}=g\cdot s_{\mathbf z}\);
a finite-dimensional subspace is closed and complemented, and
\cref{lem:hist-orthogonal-lift} gives the attained minimum. \(\square\)

\remarkheading{Nondegeneracy is a rank condition}{rem:hist-finite-tier-rank}
Nondegeneracy can be checked before any dynamics is written. With an
\(M\)-atom design \(\varrho_\ell=\sum_{d=1}^M\mathsf p_d\delta_{c_d}\), the Gram
matrix has rank at most \(MK\), so for \(n>MK\) the metric is degenerate and,
by \cref{prop:hist-semidefinite-gradient-obstruction}, the natural-gradient
equation has no solution or many. For \(\mathcal C_\ell=\R/2\pi\mathbb Z\) with
normalized arclength, \(K=1\), \(\Sigma_0=1\), \(w\equiv1\), and
\(\{\psi_j\}=\{1,\cos c,\sin c\}\), the Gram matrix is
\(\operatorname{diag}(1,\tfrac12,\tfrac12)\), with determinant \(\tfrac14\), so
\(\mathcal Q_\ell\cong\R^3\) carries a positive-definite constant metric.
\status{ESTABLISHED}

\noindent\emph{Proof.}
Under an \(M\)-atom design each field enters \(\mathsf M_\ell\) only through its
values at the atoms, so \(\mathsf M_\ell=\mathsf V^\top\mathsf D\mathsf V\), where
\(\mathsf V\in\R^{MK\times n}\) stacks the values \(\psi_j(c_d)\) and
\(\mathsf D\) is block diagonal with blocks \(\mathsf p_dw(c_d)\Sigma_0^{-1}\). On
the circle, orthogonality of
\(1,\cos c,\sin c\) under normalized arclength gives the diagonal entries
\(1,\tfrac12,\tfrac12\). \(\square\)

\propositionheading{A strong \(L^2\) tier and a weak Sobolev tier}{prop:hist-l2-sobolev-tiers}
For the integrated normal-mean metric
\(\mathsf G^{L^2}(V,W)=\int V^\top\Sigma_0^{-1}Ww d\varrho\) on
\(\mathcal Q^{L^2}=L^2(\varrho;\R^K)\), the bounds
\(0<w_-\leq w\leq w_+<\infty\) and
\(\mathrm c_-I\preceq\Sigma_0^{-1}\preceq\mathrm c_+I\) give
\(w_-\mathrm c_-\|V\|^2_{L^2}\leq\mathsf G^{L^2}(V,V)\leq w_+\mathrm c_+\|V\|^2_{L^2}\),
so the musical map is a topological isomorphism and the metric is strong; if the
\(\varrho\)-essential infimum of \(w\) is zero or its essential supremum is
infinite, the form is not equivalent to the \(L^2\) norm. On the circle
\(\mathcal C=\R/2\pi\mathbb Z\) with normalized arclength, \(K=1\), and
\(\Sigma_0=w=1\), the same metric on a Sobolev tier \(H^s\), \(s>0\), induces a
strictly coarser topology than the manifold topology, and its musical map is
injective and bounded but not surjective, so the metric is weak. The failure is
realized for \(\mathcal Q=H^1\) and \(\Fenergy(Q)=\tfrac12\int|Q'|^2\): at
\(Q=\sum_{j\geq1}j^{-2}\sin jc\), which lies in \(H^1\), a gradient would be
\(-Q''=\sum_j\sin jc\), which is not in \(L^2\), so no natural-gradient vector
field exists at \(Q\). \status{ESTABLISHED}

\noindent\emph{Proof.}
The two-sided bounds are the pointwise bounds integrated. If the essential
infimum of \(w\) is zero, the sets \(\{w<1/k\}\) have positive measure, and the
normalized indicator functions of these sets times a unit vector have
\(\mathsf G^{L^2}(V,V)\leq\mathrm c_+/k\); an infinite essential supremum gives
\(\mathsf G^{L^2}(V,V)\geq\mathrm c_-k\) in the same way. On the circle, in
Fourier coefficients \(\hat v_j\), the \(H^s\) norm is
\(\sum_j(1+j^2)^s|\hat v_j|^2\) and the metric is \(\sum_j|\hat v_j|^2\), so the two
are not equivalent. The functional
\(\ell(v)=\sum_{j\geq1}j^{s-1/2}\log(1+j)^{-1}\hat v_j\) is \(H^s\)-continuous,
since \(\sum_jj^{2s-1}\log(1+j)^{-2}(1+j^2)^{-s}<\infty\), but it is represented
by no \(L^2\) function, since \(\sum_jj^{2s-1}\log(1+j)^{-2}=\infty\). In the
example \(\|Q\|_{H^1}^2\) is a multiple of \(\sum_j(1+j^2)j^{-4}<\infty\), while
\(\|Q''\|_2^2\) is a multiple of \(\sum_j1=\infty\), and
\(d\Fenergy_Q[V]=\int Q'V'=-\int Q''V\) is \(\mathsf G^{L^2}\)-representable only
by \(-Q''\). \(\square\)

The missing ingredient on a weak Sobolev tier is representability of \(d\Fenergy_Q\)
by a vector in its declared tangent space. A bound
\(|d\Fenergy_Q[V]|\leq\mathrm c_Q\|V\|_{L^2}\) on a dense domain gives an \(L^2\)
Riesz representative, but a gradient on the \(H^s\) tier additionally requires that
representative to belong to \(H^s\). For example, on the circle with normalized
arclength, \(f(c)=\sum_{j\geq1}j^{-1}\sin(jc)\) lies in \(L^2\setminus H^1\), and
\(\Fenergy(Q)=\int fQ\) is a continuous linear functional on \(H^1\) with the displayed
bound but no \(H^1\) gradient for the \(L^2\) metric. The price of the strong \(L^2\)
tier is that \(\Fenergy\) must be \(C^1\) on \(L^2\), which excludes gradient-energy
objectives.

\section{Base averaging and the coarse configuration map}
\label{sec:hist-base-averaging}

A pointwise bundle morphism induces a map on configurations only on the
projectable set of \cref{thm:pb-section-descent}, and under the reachability
hypotheses of \cref{prop:pb-nonfunctorial-descent} that set is a proper subset
whenever the base map collapses a direction on which the vertical differential
does not vanish, and it can have empty interior in the linear tier. The coarse
configuration map \(\mathsf F_\ell\) of \eqref{eq:hist-configuration-coarse-map}
is therefore declared separately, and averaging over the fibers of a base map is
its standard instance.

\openproblemheading{Smooth structure of the projectable set}{open:hist-projectable-set}
\Cref{thm:pb-section-descent} characterizes the fine sections that descend, and
\cref{prop:pb-nonfunctorial-descent} shows that the descendable set is proper and
can have empty interior. It is not known under which hypotheses that set is a
smooth submanifold of a declared configuration manifold carrying a smooth
induced configuration map. The problem is settled by a transversality or
elliptic-regularity hypothesis on the base map and the recognition family under
which the projectable set is proved to be such a submanifold, or by a family in
which it provably is not.

\definitionheading{Averaging coarse configuration map}{def:hist-averaging-coarse-map}
Let \(f:\mathcal C_\ell\to\mathcal C_{\ell+1}\) be measurable between standard
Borel spaces, let \(\varrho\) be a finite positive measure on \(\mathcal C_\ell\)
with \(\bar\varrho=f_\#\varrho\), and let \(\{\varrho_{\bar c}\}\) disintegrate
\(\varrho\) over \(f\). Let the coarse fiber \(\bar{\mathcal B}\) be convex in a
locally convex space on which the structure group acts by restrictions of linear
maps, and let \(\Psi\) act fiberwise as the restriction of a continuous affine map
with linear part \(L\), with the integrand Bochner integrable and its barycenter
in the fiber. The averaging map is
\begin{equation}
\left(\mathsf F_\ell s\right)(\bar c)
=\int_{f^{-1}(\bar c)}\Psi\bigl(s(c)\bigr)\varrho_{\bar c}(dc).
\label{eq:hist-averaging-coarse-map}
\end{equation}
The affine chart of \(\bar{\mathcal B}\) is \emph{jointly convex} when
\((\bar\beta,u)\mapsto\bar g^F_{\bar\beta}(u,u)\) is jointly convex on
\(\bar{\mathcal B}\) times its model space; the special case in which
\(\bar g^F\) is constant in the chart is written (JC-const).

\propositionheading{Averaging reproduces descended sections}{prop:hist-averaging-descent}
In \cref{def:hist-averaging-coarse-map}, if \(\Psi\circ s\) is \(f\)-basic with
unique descended section \(\bar s\), then
\((\mathsf F_\ell s)(\bar c)=\bar s(\bar c)\) for \(\bar\varrho\)-almost every
\(\bar c\). \status{ESTABLISHED}

\noindent\emph{Proof.}
For \(\bar\varrho\)-almost every \(\bar c\) the integrand equals
\(\bar s(\bar c)\) for \(\varrho_{\bar c}\)-almost every point of the fiber, and
the disintegration is normalized. \(\square\)

\hypothesisheading{Chart joint convexity}{hyp:hist-joint-convexity}
Where \cref{thm:hist-averaging-defect} asserts a nonnegative averaging loss, the
affine chart of the coarse fiber is jointly convex in the sense of
\cref{def:hist-averaging-coarse-map}. This is a restriction on the chart, not a
property of Fisher metrics, and it excludes charts such as the covariance sector
of the Gaussian moment chart (\cref{prop:hist-averaging-convexity}).

\theoremheading{Exact averaging defect}{thm:hist-averaging-defect}
Assume \cref{def:hist-averaging-coarse-map}, \cref{hyp:hist-joint-convexity},
the fiberwise Fisher contraction
\(\Delta_F^\Psi=g^F-L^*\bar g^FL\succeq0\), which holds when \(\Psi\) is a Markov
coarse morphism by \cref{thm:pb-pullback-fisher-defect}, the weight condition
\(\bar w\circ f\leq w\) of \eqref{eq:pb-cross-scale-declarations}, and finiteness
of every integral. Then \(\mathsf F_\ell s\) is a section, \(\mathsf F_\ell\) is
gauge equivariant because the transition maps act linearly and commute with
barycenters, and for every \(Z\in T_s\mathcal Q_\ell\) the defect
\begin{equation}
\Delta_{\mathrm{avg}}(Z)
=\int_{\mathcal C_\ell}wg^F_{s(c)}(Z,Z)d\varrho
-\int_{\mathcal C_{\ell+1}}\bar w
\bar g^F_{(\mathsf F_\ell s)(\bar c)}
\bigl(T\mathsf F_\ell Z,T\mathsf F_\ell Z\bigr)d\bar\varrho
\label{eq:hist-averaging-defect}
\end{equation}
is nonnegative. Under (JC-const) it is exactly the sum of three nonnegative terms,
\begin{equation}
\Delta_{\mathrm{avg}}(Z)
=\underbrace{\int_{\mathcal C_\ell}w\Delta_F^\Psi(Z,Z)d\varrho}
_{\text{channel loss}}
+\underbrace{\int_{\mathcal C_\ell}(w-\bar w\circ f)
\bigl(L^*\bar g^F\bigr)(Z,Z)d\varrho}_{\text{weight gap}}
+\underbrace{\int_{\mathcal C_{\ell+1}}\bar w
\operatorname{Var}^{\bar g^F}_{\varrho_{\bar c}}(LZ)d\bar\varrho}
_{\text{context gap}},
\label{eq:hist-averaging-three-term}
\end{equation}
where
\(\operatorname{Var}^{\bar g}_{\varrho_{\bar c}}(U)
=\int\bar g(U,U)d\varrho_{\bar c}
-\bar g\bigl(\int Ud\varrho_{\bar c},\int Ud\varrho_{\bar c}\bigr)\).
\status{ESTABLISHED}

\noindent\emph{Proof.}
Affineness gives \(T\mathsf F_\ell Z(\bar c)=\int LZ(c)\varrho_{\bar c}(dc)\), so
\((\mathsf F_\ell s(\bar c),T\mathsf F_\ell Z(\bar c))\) is the
\(\varrho_{\bar c}\)-barycenter of \((\Psi s(c),LZ(c))\). Joint convexity and
Jensen's inequality give
\[
\bar g^F_{\mathsf F_\ell s(\bar c)}(T\mathsf F_\ell Z,T\mathsf F_\ell Z)
\leq\int\bar g^F_{\Psi s(c)}(LZ,LZ)d\varrho_{\bar c}.
\]
Multiplying by \(\bar w\),
integrating against \(\bar\varrho\), and using the disintegration turns the
right side into \(\int(\bar w\circ f)(L^*\bar g^F)(Z,Z)d\varrho\); the weight
condition and \(\Delta_F^\Psi\succeq0\) then give nonnegativity. When
\(\bar g^F\) is constant in the chart, the Jensen step is an equality minus the
variance term, and expanding \(g^F=\Delta_F^\Psi+L^*\bar g^F\) separates the
channel and weight contributions, which gives
\eqref{eq:hist-averaging-three-term}. \(\square\)

\Cref{thm:hist-pointwise-contraction-lift} is the case of a shared base, where
the disintegration is a point mass and no convexity is needed; the averaging
defect admits a collapsing base map at the price of affine-fiber, convexity, and
barycenter hypotheses, so no general comparison of logical strength between the
two results is asserted.

\propositionheading{Joint convexity cannot be dropped}{prop:hist-averaging-convexity}
The joint-convexity hypothesis holds in the location sector of a Gaussian fiber
with fixed covariance, where the metric is constant, and in the law chart, where
the integrand \((p,\dot p)\mapsto\dot p^2/p\) has Hessian determinant zero and
positive trace. It fails in the covariance sector of the Gaussian moment chart,
where \((\Sigma,\dot\Sigma)\mapsto\dot\Sigma^2/(2\Sigma^2)\) has Hessian
determinant \(-\dot\Sigma^2\Sigma^{-6}<0\) for \(\dot\Sigma\ne0\), and there the
averaging defect can be negative: for two contexts with the uniform
disintegration, unit weights, \(\Psi\) the identity, the centered normal fiber
\(\{\mathcal N(0,\Sigma)\}\), the configuration \((\Sigma_1,\Sigma_2)=(1,\delta)\),
and the tangent \(Z=(1,0)\), the fine integrated metric is \(\tfrac14\), the
image metric is \(1/(2(1+\delta)^2)\), their ratio is \(2/(1+\delta)^2\), equal to
\(\tfrac{20000}{10201}\) at \(\delta=10^{-2}\), and the defect is negative exactly
for \(\delta<\sqrt2-1\), tending to \(-\tfrac14\) as \(\delta\to0^+\). Every other
hypothesis of \cref{thm:hist-averaging-defect} holds on this datum.
\status{ESTABLISHED}

\noindent\emph{Proof.}
The Hessians are direct computations. With the uniform disintegration the fine
metric is \(\tfrac12\cdot\tfrac12+\tfrac12\cdot0=\tfrac14\); the barycenter
configuration is \(\bar\Sigma=(1+\delta)/2\) with tangent \(\tfrac12\), whose
metric is \((\tfrac12)^2/(2\bar\Sigma^2)=1/(2(1+\delta)^2)\). The defect
\(\tfrac14-1/(2(1+\delta)^2)\) is negative exactly when \((1+\delta)^2<2\).
\(\Delta_F^\Psi=0\) and the weights are equal, so only joint convexity fails.
\(\square\)

No contraction statement may therefore be attached to a generic averaging or
variational coarse map.

\propositionheading{Smoothness and gauge descent at the exhibited tier}{prop:hist-coarse-map-smoothness}
Take the averaging map \eqref{eq:hist-averaging-coarse-map} with
\(\bar\varrho=f_\#\varrho\), its disintegration \(\{\varrho_{\bar c}\}\), and
\(\Psi\) acting on means as the identity. Extend
\cref{def:hist-finite-configuration-tier} to scale \(\ell+1\) with base
\(\mathcal C_{\ell+1}\), measure \(\bar\varrho\), weight \(\bar w\), coarse fiber
\(\{\mathcal N(\mu,\bar\Sigma_0)\}\), and a coarse basis
\(\bar\psi_1,\dots,\bar\psi_{\bar n}\) linearly independent in
\(L^2(\bar w\bar\varrho;\bar\Sigma_0^{-1})\) whose span contains every constant
field, with constant coefficient map \(\jmath_{\ell+1}\) as in
\eqref{eq:hist-finite-tier-coefficient-action}. Suppose the fiberwise average of
each fine basis field lies in the coarse span,
\(\int_{f^{-1}(\bar c)}\psi_jd\varrho_{\bar c}=\sum_k\mathbf E_{kj}\bar\psi_k(\bar c)\)
for \(\bar\varrho\)-almost every \(\bar c\), with a constant matrix
\(\mathbf E\in\R^{\bar n\times n}\). Then the averaging map reads
\(\mathsf F_\ell(\mathbf z)=\mathbf E\mathbf z\) in coordinates; it is linear,
hence \(C^\infty\), with tangent map \(\mathbf E\); it intertwines the coefficient
representations of the constants,
\begin{equation}
\mathbf E\jmath_\ell=\jmath_{\ell+1},
\label{eq:hist-finite-tier-coefficient-intertwining}
\end{equation}
so it is gauge equivariant, \(\mathbf E\mathscr A^\ell_g=\mathscr A^{\ell+1}_g\mathbf E\);
and its metric compatibility \(\mathsf F_\ell^*\mathsf G_{\ell+1}\preceq\mathsf G_\ell\)
is the matrix inequality
\begin{equation}
\mathbf E^\top\overline{\mathsf M}\mathbf E\preceq\mathsf M_\ell,
\qquad
\overline{\mathsf M}_{kk'}
=\int\bar\psi_k^\top\bar\Sigma_0^{-1}\bar\psi_{k'}\bar wd\bar\varrho,
\label{eq:hist-coarse-map-compatibility}
\end{equation}
which holds whenever \(\bar w\circ f\leq w\) \(\varrho\)-almost everywhere and
\(\bar\Sigma_0\succeq\Sigma_0\). Neither sufficient condition can be dropped:
with \(f\) the identity, \(\bar\psi=\psi\), and \(\mathbf E=I\), the choice
\(\bar\Sigma_0=\Sigma_0\), \(\bar w=100w\) gives
\(\mathbf E^\top\overline{\mathsf M}\mathbf E-\mathsf M_\ell=99\mathsf M_\ell\succ0\),
and the choice \(\bar w=w\), \(\bar\Sigma_0=\Sigma_0/2\) gives
\(\mathbf E^\top\overline{\mathsf M}\mathbf E-\mathsf M_\ell=\mathsf M_\ell\succ0\).
\status{ESTABLISHED}

\noindent\emph{Proof.}
Averaging a configuration of the declared family produces the mean field
\[
\bar c\longmapsto\sum_jz_j\int\psi_jd\varrho_{\bar c}
=\sum_k(\mathbf E\mathbf z)_k\bar\psi_k(\bar c),
\]
which is the coarse configuration \(s_{\mathbf E\mathbf z}\). Averaging the
identity \(\sum_j(\jmath_\ell g)_j\psi_j(c)=g\) against the probability measure
\(\varrho_{\bar c}\) gives
\[
\sum_k(\mathbf E\jmath_\ell g)_k\bar\psi_k(\bar c)=g=\sum_k(\jmath_{\ell+1}g)_k\bar\psi_k(\bar c)
\]
almost everywhere, and linear independence of the \(\bar\psi_k\) gives
\eqref{eq:hist-finite-tier-coefficient-intertwining}; then
\(\mathbf E\mathscr A^\ell_g(\mathbf z)=\mathbf E\mathbf z+\mathbf E\jmath_\ell g
=\mathscr A^{\ell+1}_g(\mathbf E\mathbf z)\). Both metrics are constant Gram
forms, so the pullback of \(\mathsf G_{\ell+1}\) is
\(\mathbf E^\top\overline{\mathsf M}\mathbf E\). For \(\mathbf v\in\R^n\) and
\(U=\sum_jv_j\psi_j\), the two conditions give
\(wU^\top\Sigma_0^{-1}U\geq(\bar w\circ f)U^\top\bar\Sigma_0^{-1}U\) pointwise,
and the disintegration followed by Jensen's inequality for the convex quadratic
form \(\bar\Sigma_0^{-1}\) gives
\[
\mathbf v^\top\mathsf M_\ell\mathbf v
\geq\int\bar w\int U^\top\bar\Sigma_0^{-1}Ud\varrho_{\bar c}d\bar\varrho
\geq\int\bar w
\Bigl(\int Ud\varrho_{\bar c}\Bigr)^\top\bar\Sigma_0^{-1}
\Bigl(\int Ud\varrho_{\bar c}\Bigr)d\bar\varrho
=\mathbf v^\top\mathbf E^\top\overline{\mathsf M}\mathbf E\mathbf v.
\]
This is the (JC-const) case of \cref{thm:hist-averaging-defect}, with channel loss
\(\Sigma_0^{-1}-\bar\Sigma_0^{-1}\). In the two counterexamples
\(\overline{\mathsf M}=100\mathsf M_\ell\) and \(\overline{\mathsf M}=2\mathsf M_\ell\);
for the Gram matrix \(\operatorname{diag}(1,\tfrac12,\tfrac12)\) of
\cref{rem:hist-finite-tier-rank} the first difference has smallest eigenvalue
\(49.5\). \(\square\)

\propositionheading{Integrated horizontal conformality}{prop:hist-integrated-conformality}
Let the two scales share the base \(\mathcal C_i\), the finite base measure
\(\varrho_i\), and the positive channel weights \(w_x\); let
\(\mathsf F_\ell\) act pointwise as \((\mathsf F_\ell Q)(c)=\Psi_c(Q(c))\), each
\(\Psi_c\) the pushforward of a parameter-independent Markov kernel; let both
configuration metrics be the weighted integrals of the fiber Fisher metrics; and
let every integral be finite. Then for \(Z,W\in T_Q\mathcal Q_\ell\)
\[
\mathsf G_\ell(Z,W)-\mathsf G_{\ell+1}(T\mathsf F_\ell Z,T\mathsf F_\ell W)
=\int_{\mathcal C_i}\sum_xw_x\Delta_F^{\Psi_c}\bigl(Z_x(c),W_x(c)\bigr)d\varrho_i(c),
\]
the integrated form of \eqref{eq:hist-pointwise-contraction-lift}. If
\(\mathcal H=(\ker T\mathsf F_\ell)^{\perp_{\mathsf G_\ell}}\) and
\(T\mathsf F_\ell|_{\mathcal H}\) is onto, as for the submersions of
\cref{prop:hist-natural-gradient-sufficiency}, then \(\mathsf F_\ell\) is
horizontally conformal with dilation \(\varphi_\ell\) exactly when
\begin{equation}
\int_{\mathcal C_i}\sum_xw_x
\Delta_F^{\Psi_c}\bigl(Z_x(c),W_x(c)\bigr)d\varrho_i(c)
=(1-\varphi_\ell^2)\mathsf G_\ell(Z,W)
\quad\text{for every }Z,W\in\mathcal H.
\label{eq:hist-integrated-conformality-criterion}
\end{equation}
The integrated identity does not imply the pointwise one: on a base consisting
of two contexts of weight \(\tfrac12\), with fine fibers \(\mathcal N(\mu,1)\),
the constant-mean tangent, and independent Gaussian noises of variances \(3\) and
\(\tfrac13\) added at the two contexts, the pulled-back location Fisher forms are
\(\tfrac14\) and \(\tfrac34\), the integrated ratio is \(\varphi_\ell^2=\tfrac12\),
and the pointwise ratios differ. If \(\mathcal H\) is closed under multiplication
by nonnegative compactly supported smooth functions on the base, the integrated
identity implies the pointwise identity \(\varrho_i\)-almost everywhere.
\status{ESTABLISHED}

\noindent\emph{Proof.}
Differentiating the pointwise action and inserting it into the coarse metric
gives the first display, whose integrand is the polarized vertical Fisher
defect \eqref{eq:pb-vertical-fisher-defect} of \(\Psi_c\). Horizontal
conformality means
\(\mathsf G_{\ell+1}(T\mathsf F_\ell Z,T\mathsf F_\ell W)=\varphi_\ell^2\mathsf G_\ell(Z,W)\)
on \(\mathcal H\), with \(T\mathsf F_\ell|_{\mathcal H}\) onto; subtracting from
\(\mathsf G_\ell(Z,W)\) and using the first display gives the equivalence. In the
example the pushforward of \(\mathcal N(\mu,1)\) by noise of variance \(\sigma^2\)
is \(\mathcal N(\mu,1+\sigma^2)\), whose location Fisher information is
\(1/(1+\sigma^2)\), namely \(\tfrac14\) and \(\tfrac34\); the fine metric of the
constant-mean tangent is \(1\) and the coarse metric is
\(\tfrac12(\tfrac14+\tfrac34)=\tfrac12\). For the last claim, apply the
integrated identity to \(\chi Z\in\mathcal H\) for smooth compactly supported
\(\chi\geq0\): the integrable function
\(c\mapsto\sum_xw_x[\Delta_F^{\Psi_c}-(1-\varphi_\ell^2)g^F_x](Z_x(c),Z_x(c))\)
integrates to zero against every \(\chi^2\); monotone limits of such \(\chi^2\)
give the indicator functions of compact coordinate balls, so its integral over
every such ball vanishes and the function vanishes almost everywhere.
Polarization gives the bilinear identity. \(\square\)

\section{Proof of the strong Hilbert quotient theorem}
\label{sec:hist-strong-hilbert-quotient-proof}

\noindent\emph{Proof of \cref{thm:hist-strong-hilbert-quotient}.}
Closedness of \(\mathfrak a\) with compact fibers implies compact inverse images
of compact sets: an open cover of such an inverse image has a finite subcover on
each fiber, closedness gives a neighborhood of each target point whose inverse
image lies in the union of that subcover, and finitely many such neighborhoods
cover the compact target set. Conversely, compact-preimage properness gives
closedness, since a convergent sequence of images of a closed set lifts into the
compact inverse image of the sequence and its limit, and first countability
gives a convergent subsequence of lifts.

The length distance of a strong metric induces the manifold topology on each
connected component. In a Hilbert chart at \(Q\), strongness and operator-norm
continuity give constants \(0<\mathrm c_1<\mathrm c_2<\infty\) and a coordinate
ball on which
\begin{equation}
\mathrm c_1\|u\|^2\leq\mathsf g_y(u,u)\leq\mathrm c_2\|u\|^2.
\label{eq:hist-strong-chart-bounds}
\end{equation}
Straight segments give the upper distance bound, and a path from the center that
leaves a smaller ball of radius \(\delta_0\) has length at least
\(\sqrt{\mathrm c_1}\delta_0\),
so small length distance forces the endpoint into every prescribed
neighborhood; no minimizing geodesic is used.

For \(c_\xi(t)=\exp(t\xi)\cdot Q\) and \(\mathscr A_g(Q)=g\cdot Q\), the group law
and isometry give
\begin{equation}
\dot c_\xi(t)=T_Q\mathscr A_{\exp(t\xi)}A_Q\xi,
\qquad
\|\dot c_\xi(t)\|_{\mathsf g}=\|A_Q\xi\|_{\mathsf g_Q},
\qquad
d_{\mathsf g}\bigl(Q,\exp(\xi)\cdot Q\bigr)\leq\|A_Q\xi\|_{\mathsf g_Q}.
\label{eq:hist-orbit-constant-speed}
\end{equation}
If \(A_Q\xi=0\), the curve is constant, freeness gives \(\exp(t\xi)=e\) for
every \(t\), and hence \(\xi=0\).

Suppose that unit vectors \(\xi_n\) in the Banach norm of \(\mathfrak g\) satisfied
\(\|A_Q\xi_n\|_{\mathsf g_Q}\to0\). Choose \(\epsilon>0\) whose closed ball in
\(\mathfrak g\) lies in an exponential chart, and set \(g_n=\exp(\epsilon\xi_n)\).
Then \eqref{eq:hist-orbit-constant-speed} and the metric topology give
\(g_n\cdot Q\to Q\). Properness over the compact set
\(\{(g_n\cdot Q,Q)\}\cup\{(Q,Q)\}\) gives a subsequence \((g_n,Q)\to(g,Q)\) with
\(g\cdot Q=Q\), freeness gives \(g=e\), and the continuous local logarithm gives
\(\epsilon\xi_n\to0\), a contradiction. Therefore, for some
\(\mathrm c_Q>0\),
\begin{equation}
\|A_Q\xi\|_{\mathsf g_Q}\geq\mathrm c_Q\|\xi\|_{\mathfrak g}
\quad\text{for every }\xi\in\mathfrak g.
\label{eq:hist-orbit-coercivity}
\end{equation}
This is vacuous if \(\mathfrak g=\{0\}\); otherwise a convergent sequence
\(A_Q\xi_n\) makes \(\xi_n\) Cauchy, and completeness puts its limit in the range,
which is closed, in agreement with \citet[Theorem~2.1]{JotzNeeb2008}.

Fix \(Q_0\). The inner product
\(\langle\xi,\eta\rangle_{\mathfrak g,H}=\mathsf g_{Q_0}(A_{Q_0}\xi,A_{Q_0}\eta)\)
has a complete norm equivalent to the Banach norm by
\eqref{eq:hist-orbit-coercivity}, so \(\mathfrak g\) is Hilbertizable; adjoints
are taken for this fixed Hilbert norm. In a tangent chart write
\(\mathsf g_y(v,u)=\langle v,M_yu\rangle\) and define
\begin{equation}
B_y=A_y^*M_yA_y,
\qquad
P_y=A_yB_y^{-1}A_y^*M_y,
\qquad
\xi_*(y,v)=B_y^{-1}A_y^*M_yv.
\label{eq:hist-smooth-orbit-projection}
\end{equation}
The action makes \(A_y\) smooth in operator norm. Near each point,
\eqref{eq:hist-orbit-coercivity} and continuity of the metric give
\(B_y\geq\mathrm c_0I\) for one \(\mathrm c_0>0\); such a self-adjoint operator
has closed range, since \(\|B_y\xi\|\geq\mathrm c_0\|\xi\|\), and dense range,
since its orthogonal complement is \(\ker B_y=\{0\}\), so it is invertible, and
the local Neumann series gives smoothness of the inverse.

Direct multiplication gives \(P_y^2=P_y\), \(\operatorname{ran}P_y=V_y\), and
\(A_y^*M_y(v-P_yv)=0\), so \(P_y\) is the metric-orthogonal vertical projection.
Its ranges and kernels are smooth subbundles: near a point \(y_0\) the smooth
operator
\begin{equation}
S_y=P_yP_{y_0}+(I-P_y)(I-P_{y_0}),
\qquad
P_yS_y=S_yP_{y_0},
\label{eq:hist-projection-trivialization}
\end{equation}
equals \(I\) at \(y_0\) and remains invertible nearby, which identifies the
splitting at \(y_0\) with that at \(y\). Orthogonality gives
\begin{equation}
\|v-A_y\xi\|_{\mathsf g_y}^2
=\|(I-P_y)v\|_{\mathsf g_y}^2
+\|A_y(\xi-\xi_*(y,v))\|_{\mathsf g_y}^2,
\label{eq:hist-quotient-pythagoras}
\end{equation}
which proves attainment and uniqueness of the minimizer before any quotient
chart is used.

For the quotient charts, properness and freeness imply that for each identity
neighborhood \(\mathcal W\subseteq\mathbb G\) there is a neighborhood
\(\mathcal O\) of \(Q\) with
\begin{equation}
g\mathcal O\cap\mathcal O\ne\varnothing\quad\Longrightarrow\quad g\in\mathcal W.
\label{eq:hist-proper-local-return}
\end{equation}
Otherwise there are \(y_n,z_n\to Q\) and \(g_n\notin\mathcal W\) with
\(g_ny_n=z_n\); properness applied to the compact set of pairs \((z_n,y_n)\) and
its limit gives a convergent subsequence of \((g_n,y_n)\) whose limiting group
element fixes \(Q\), hence is \(e\), a contradiction. Choose a local submanifold
\(\mathcal S_0\) through \(Q\), obtained in a chart from the closed complemented
space \(H_Q\). The map \((g,s)\mapsto g\cdot s\) on \(\mathbb G\times\mathcal S_0\)
has derivative \((\xi,u)\mapsto A_Q\xi+u\) at \((e,Q)\), a bounded isomorphism by
the orthogonal splitting and \eqref{eq:hist-orbit-coercivity}, so the Banach
inverse function theorem gives a product \(\mathcal W\times\mathcal S_1\) on which
it is a diffeomorphism onto an open set. Shrink to
\(\mathcal S\subseteq\mathcal S_1\cap\mathcal O\) using
\eqref{eq:hist-proper-local-return} for this \(\mathcal W\). If \(g\cdot s=t\) for
\(s,t\in\mathcal S\), then \(g\in\mathcal W\), and injectivity of the local product
gives \(g=e\) and \(s=t\). Consequently
\begin{equation}
\mathbb G\times\mathcal S\longrightarrow\mathbb G\cdot\mathcal S,
\qquad
(g,s)\longmapsto g\cdot s,
\label{eq:hist-principal-slice}
\end{equation}
is bijective onto an open set, being a union of translates of the local product,
and a local diffeomorphism at every point by translation; it is a
diffeomorphism. These are principal trivializations, with smooth transition maps
given by the slice components of their inverses.

The orbit relation is closed by the closedness of \eqref{eq:hist-proper-action-map},
and \(\pi\) is open because saturations of open sets are open, so a product
neighborhood disjoint from the orbit relation separates two distinct orbits, and
the quotient is Hausdorff. Truncating the strong length distance at \(1\), and
giving points in different components distance \(1\), gives an invariant metric
\(d_0\) inducing the topology of \(\mathcal M\), and
\begin{equation}
\bar d\bigl([Q],[R]\bigr)=\inf_{g\in\mathbb G}d_0(Q,g\cdot R)
\label{eq:hist-quotient-topology-metric}
\end{equation}
is positive for distinct orbits because every orbit is closed; symmetry follows
by inversion, and the triangle inequality by isometry and multiplication in
\(\mathbb G\). Its balls give the quotient topology, since their inverse images
are open and an open saturated neighborhood of \(Q\) contains a \(d_0\)-ball at
\(Q\). The quotient is therefore metrizable, hence paracompact, as the standard
definitions of Hilbert manifolds require; \(\bar d\) is not asserted to be the
Riemannian quotient distance.

The action transports \(V_Q\) by
\(T\mathscr A_gA_Q\xi=A_{gQ}\operatorname{Ad}_g\xi\), and isometry transports
\(H_Q\) as well. Define
\begin{equation}
\overline{\mathsf g}_{[Q]}(d\pi_Qv,d\pi_Qu)
=\mathsf g_Q\bigl((I-P_Q)v,(I-P_Q)u\bigr).
\label{eq:hist-strong-quotient-metric}
\end{equation}
Adding vertical vectors to the lifts leaves this unchanged, and equivariance of
the orthogonal projection gives independence of \(Q\) within its orbit. In each
slice it is smooth, and \((I-P_s)|_{T_s\mathcal S}:T_s\mathcal S\to H_s\) is a
bounded isomorphism because the local product gives
\(T_s\mathcal M=V_s\oplus T_s\mathcal S\), so the pulled-back inner product is
strong. Equations \eqref{eq:hist-quotient-pythagoras} and
\eqref{eq:hist-strong-quotient-metric} give \eqref{eq:hist-quotient-gauge-speed}.

Finally, for \(\widetilde Q(t)=g(t)\cdot Q(t)\),
\begin{equation}
\dot{\widetilde Q}(t)
=T\mathscr A_{g(t)}\dot Q(t)+A_{\widetilde Q(t)}\delta^rg(t),
\label{eq:hist-moving-quotient-representative}
\end{equation}
where \(\delta^rg\) is the right logarithmic derivative of \(g\). The second term
is vertical and the first preserves the horizontal norm, which gives
representative independence of speed and duration. The unique coefficient
\(\xi_*(Q,v)\) also defines a smooth principal connection, since it sends
\(A_Q\xi\) to \(\xi\) and transforms by \(\operatorname{Ad}_g\) by uniqueness and
equivariance of \(P\); no global gauge section is needed. \(\square\)
